\chapter{Lelong number}

\begin{thebibliography}{n}
\bibitem[Dem12]{Dem}
J.-P. Demailly,
\textit{Complex Analytic and Differential Geometry},
Book, 2012.
\url{https://people.math.harvard.edu/~demarco/Math274/Demailly_ComplexAnalyticDiffGeom.pdf}\footnote{Demailly先生のページがどうも開けず, Demarco先生のページのリンクを引用した. Demailly先生のページは一回見れないことがあったので, もしかしたらそのうちなくなる可能性もある. なおそれを見越して私はDemailly先生にあったpdfを全てダウンロードした.}
\bibitem[Rud]{Rud}
W. Rudin. \textit{Functional analysis.} 2nd edn. International Series in Pure and Applied Mathematics. McGraw-Hill, Inc., New York. (1991.)
\bibitem[Rud 2]{Rud2}
W. Rudin. \textit{Real and complex analysis. } 3rd ed. New York, NY: McGraw-Hill 
\bibitem[NO]{NO}
J. Noguchi, T.Ochiai \textit{Geometric Function Theory in Several Complex Variables} Translations of Mathematical Monographs
Volume: 80; 1990; 282 pp
\end{thebibliography}
 
カレントに関しては\cite[Chapter 3]{NO}を参考にしている. 
またpositive currentに関しては, 
\cite[Chapter 1. Section 3.C]{Dem}, \cite[Chapter 3. Section 1.]{Dem}
などを参考にしている. 

% [page 1]
\section{Lelong number}

\underline{Notation}
\begin{itemize}
\item $(z^1,\ldots,z^m)$: coordinate of $\mathbb{C}^m$
\item $d^c:=\dfrac{1}{2\pi}\dfrac{i}{2}(\bar{\partial}-\partial)$ \quad (thus $dd^c=\dfrac{i}{2\pi}\partial\bar{\partial}$)
\item $\alpha:=dd^c|z|^2=\dfrac{i}{2\pi}\displaystyle\sum_{j=1}^m dz^j\wedge d\bar{z}^j$
\item $\beta:=dd^c\log|z|^2=\dfrac{i}{2\pi}\displaystyle\sum_{j,k=1}^m\left(\dfrac{\delta_{jk}}{|z|^2}-\dfrac{z^k\bar{z}^j}{|z|^4}\right)dz^j\wedge d\bar{z}^k$
\item $\alpha^k,\beta^k$：$k$ 回ウェッジ積，$\alpha^0:=1$，$\beta^0:=1$．
\item $\eta:=d^c\log|z|^2\wedge\beta^{m-1}$
\item $B(z;r):=\{w\in\mathbb{C}^m\mid |w-z|<r\}$
\item $\Gamma(z;r):=\{w\in\mathbb{C}^m\mid |w-z|=r\}$
\item $B(r):=B(0;r)$，$\Gamma(r):=\Gamma(0;r)$．
\end{itemize}

\begin{rem}
\label{rem-H-5.1}
%\underline{Rem 1}
\begin{equation*}
d\alpha=d\beta=0,\qquad
\alpha^m=\frac{m!}{\pi^m}\bigwedge_{j=1}^m\frac{i}{2}dz^j\wedge d\bar{z}^j,\qquad
d\eta=\beta^m=0.
\end{equation*}

$\beta^m=0$ については
$A:=\left(\dfrac{\delta_{jk}}{|z|^2}-\dfrac{z^k\bar{z}^j}{|z|^4}\right)_{\substack{1\leq j\leq m\\1\leq k\leq m}}$
として $\det(A)=0$ を示せばよいが：
\begin{equation*}
|z|^4A
=
|z|^2E_m
-
\begin{pmatrix}
\bar{z}^1\\
\vdots\\
\bar{z}^m
\end{pmatrix}
\begin{pmatrix}
z^1&\cdots&z^m
\end{pmatrix}
\end{equation*}
より
$A\begin{pmatrix}\bar{z}^1\\ \vdots\\ \bar{z}^m\end{pmatrix}=0$
なので行列式の性質から$\det(A)=0$．
\end{rem}

% [page 2]
%\underline{Fact 2}
\begin{rem}
\label{rem-H-5.2}
半径 $r$ の $d$ 次元球の体積はガンマ関数を用いて, 
$\frac{\pi^{d/2}}{\Gamma\left(\frac{d}{2}+1\right)}r^d
$となる. よって$dz^j=dx^j+i\,dy^j$と$d=2m$の代入から
\begin{equation*}
\int_{B(z;r)}\alpha^m
=
\frac{m!}{\pi^m}
\int_{B(z;r)}
dx^1\wedge dy^1\wedge\cdots\wedge dx^m\wedge dy^m
=
r^{2m}
\end{equation*}
ここに$\Gamma\left(\frac{d}{2}+1\right)=\Gamma\left(m+1\right)=m!$を用いた
\end{rem}

\begin{tcolorbox}[mybox]
\begin{lem}
\label{lem-H-5.3}
%\underline{Lem 3}
\begin{equation*}
\left.\beta\right|_{\Gamma(r)}
=
\frac{1}{r^2}\left.\alpha\right|_{\Gamma(r)}.
\end{equation*}
\end{lem}
\end{tcolorbox}
\begin{proof}
$f\colon\mathbb{C}^m\to\mathbb{R}$を$f(z):=|z|^2$とする. 
すると$\Gamma(r)$の定義と$f$の定義式から
\begin{equation*}
\left.(\partial f)\right|_{\Gamma(r)}
+
\left.(\bar{\partial}f)\right|_{\Gamma(r)}
=
\left.(df)\right|_{\Gamma(r)}
=
0.
\end{equation*}
よって
\begin{align*}
\beta
&\underalign{\text{$\beta$の定義}}{=}
\frac{i}{2\pi}\partial\bar{\partial}\log f(z)
=
\frac{i}{2\pi}\partial\frac{\bar{\partial}f(z)}{|z|^2}
=
\frac{i}{2\pi}
\left(
\frac{\partial\bar{\partial}f(z)}{|z|^2}
+
\frac{\bar{\partial}f(z)\wedge\partial f(z)}{|z|^4}
\right)
%\\
%&=
%\frac{1}{|z|^2}\alpha
%+
%\frac{i}{2\pi}\cdot
%\frac{\bar{\partial}f(z)\wedge\partial f(z)}{|z|^4}.
\end{align*}
以上より
\begin{equation*}
\left.\beta\right|_{\Gamma(r)}
=
\frac{1}{r^2}\left.\alpha\right|_{\Gamma(r)}
+
\frac{i}{2\pi}\cdot
\frac{
\overbrace{-\partial f(z)}^{\bar{\partial}f(z)}
\wedge\partial f(z)
}{r^4}
=
\frac{1}{r^2}\left.\alpha\right|_{\Gamma(r)}.
\end{equation*}
\end{proof}

\begin{rem}[Poincare Lemma]
\label{rem-H-5.4}
%\underline{Fact 4 (Poincaré lemma)}
$U\subset\mathbb{R}^d$を星型領域とする. 
任意の$\varphi\in\mathcal{E}^p(U)$について, $\psi\in\mathcal{E}^{p-1}(U)$があって, $d\psi=\varphi$.
\end{rem}

% [page 3]
%\underline{Def 5 (Radon measure, [松澤 2026-04-13, Def 7 以降])}
\begin{tcolorbox}[mybox]
\begin{defn}
$U\subset\mathbb{R}^d$を開集合, 
$\mathcal{B}(U)$を全ての開集合を含む最小の Borel $\sigma$代数とする. 

Radon measure' $\mu\colon\mathcal{B}(U)\to[0,\infty]$はBorel measureであり, 
\begin{enumerate}
\item (正則性) 任意の$ E\in\mathcal{B}(U)$について, 
\begin{equation*}
\mu(E)
=
\inf\{\mu(V)\mid E\subset V\subset U;\ \text{open}\}
=
\sup\{\mu(A)\mid A\subset E:\ \text{cpt}\}.
\end{equation*}
\item 任意のコンパクト集合$K\in\mathcal{B}(U)$について$\mu(K)<\infty$.
\end{enumerate}

complex Radon measure $\mu\colon\mathcal{B}(U)\to\mathbb{C}$を複素Borel測度でtotal variation
\begin{equation*}
\mathcal{B}(U)\longrightarrow[0,\infty] \quad E \longmapsto
|\mu|(E)
:=
\sup\left\{
\sum_{i=1}^{\infty}|\mu(E_i)| \mid E=\bigcup_{i=1}^{\infty}E_i, E_i\cap E_j=\varnothing (\forall i,j)
\right\}.
\end{equation*}
%$|\mu|\colon\mathcal{B}(U)\to[0,\infty]$
がRadon測度となることとする.
\end{defn}
\end{tcolorbox}
上はRudinの意味でのRadon measureと同じである. 
完全加法性が絶対収束になることを要請しており, この場合total variationが有限になること従う. 

%別の流儀では cpx. Radon meas. を $\mathcal{B}(U)$ ではなく$\{E\in\mathcal{B}(U)\mid\overline{E}:\text{cpt}\}$の枠組みで定義している．このときは$E\in\mathcal{B}(U):\text{non-rel. cpt.}$に対して $\mu(E)$ は定義されない．


%\underline{Obs 8}
\begin{rem}
\label{rem-H-5.8}
$U\subset\mathbb{C}^m$ open,  $T\in\mathcal{D}'^{(p,p)}(U)$ positive currentとする. 
$\alpha$をstrongly positive formとすると, $T\wedge\alpha^{m-p}$もpositiveである. 特に$T\wedge\alpha^{m-p}\in\mathcal{D}'^{(m,m)}(U)$ となる. 
よってRieszの表現定理から$\exists!\mu_{T\wedge\alpha^{m-p}}$がpositive real Radon measureとして存在し次を満たす. 
\begin{equation*}
(T\wedge\alpha^{m-p})(f)
=
\int_U f\,d\mu_{T\wedge\alpha^{m-p}}
\qquad
(\forall f\in\mathcal{K}(U)).
\end{equation*}
\end{rem}

\begin{tcolorbox}[mybox]
\begin{defn}
\label{defn-H-5.9}
%\underline{Def 9}
$U, T$を上のようにとる. 
任意の$z\in U$, $r\in\mathbb{R}_{>0}$で$\overline{B}(z;r)\subset U$となるものについて, 
\begin{equation*}
n(z;r,T)
:=
\frac{1}{r^{2(m-p)}}
\mu_{T\wedge\alpha^{m-p}}\bigl(B(z;r)\bigr),
\end{equation*}
%where, for $\forall S\subset\mathbb{C}^m$,
%\begin{equation*}
%\chi_S(z)
%:=
%\begin{cases}
%1 & (z\in S),\\
%0 & (z\notin S).
%\end{cases}
%\end{equation*}
とする. 
ここで$\mu_{T\wedge\alpha^{m-p}}$がpositive measureなので, $n(z;r,T)\geq 0$である. 
\end{defn}
\end{tcolorbox}


\begin{tcolorbox}[mybox]
\begin{thm}
\label{thm-H-5.10}
%\underline{Thm 10}
$U, T$を上のようにとる. さらに$dT=0$を仮定する. 
この時$z\in U$について, $r\longmapsto n(z;r,T)$は単調非減少関数である.
\end{thm}
\end{tcolorbox}

% [page 1]
証明には2つの道具が要る.

\begin{tcolorbox}[mybox]
\begin{lem}[{Proposition \ref{prop-M-8.4.9}}]
\label{lem-H-5.11}
%\noindent\underline{\textbf{Lem 11 (On trace)}}
$U\subset\mathbb{C}^m$ openとするとき, 
\begin{equation*}
T=i^{p^2}\sum_{J,K\in\{m;p\}}T_{J\bar K}\,dz^J\wedge d\bar z^K
\in\mathcal{D}'^{(p,p)}(U)
\end{equation*}
は positive currentであり, 
\begin{equation*}
\operatorname{Tr}(T):=\sum_{J\in\{m;p\}}T_{J\bar J}
\end{equation*}
は positive $(0,0)$-currentである. 
\end{lem}
\end{tcolorbox}

\begin{tcolorbox}[mybox]
\begin{defn}
\label{defn-H-5.12}
%\noindent\underline{\textbf{Def 12 (Smoothing)}}
\begin{itemize}
\item $h\colon\mathbb{C}^m\longrightarrow\mathbb{R}$を $C^\infty$関数で$h(z)\geq 0$かつ$\operatorname{Supp}(h)\subset B(0;1)$となるもので, 
$h(z)$は$\| z\|$にしかよらず
\begin{equation*}
\int_{\mathbb{R}^m}h(z)\,dz=1.
\end{equation*}
となるもの. 
\item $h_\epsilon\colon\mathbb{C}^m\longrightarrow\mathbb{R}$を $C^\infty$関数で
\begin{equation*}
h_\epsilon(z):=\frac{1}{\epsilon^m}h\left(\frac{z}{\epsilon}\right).
\end{equation*}
\item 
$
T=\sum_{\substack{J\in\{m;p\}\\K\in\{m;q\}}}
T_{J\bar K}\,dz^J\wedge d\bar z^K
\in\mathcal{D}'^{(p,q)}(U)
$
current を開集合$U\subset\mathbb{C}^m $のcurrent. 
ここで$T_{J\bar K}$は
$$T_{J\bar K}\wedge
\left(\bigwedge_{\ell=1}^m dz^\ell\wedge d\bar z^\ell\right)
=
T\wedge dz^{J^\circ}\wedge d\bar z^{K^\circ}
$$
で定義する. 
$J^\circ,K^\circ$は$J, K$のordered complementである. 
\item 
$
U_\epsilon:=
\left\{
z\in\mathbb{C}^m
\ \middle|\
\overline{B(z;\epsilon)}\subset U
\right\}.
$
\end{itemize}
この時, 
\begin{equation*}
T_\epsilon:=
\sum_{\substack{J\in\{m;p\}\\K\in\{m;q\}}}
(T_{J\bar K}*h_\epsilon)\,dz^J\wedge d\bar z^K
\end{equation*}
 は $\mathcal{D}'^{(p,q)}(U_\epsilon)$の元である. 
 ここで
\begin{equation*}
T_{J\bar K}*h_\epsilon(z):= T_{J\bar K}\bigl(h_\epsilon(\bullet-z)\bigr)
\end{equation*}
とする. ここで$h_\epsilon(\bullet-z)$は変数$\bullet$とする関数なので, $T_{J\bar K}$を当てることができる. 

\begin{equation*}
(T_{J\bar K}*h_\epsilon)(f)
:=
\int_{U_\epsilon}
\left(T_{J\bar K}\bigl(h_\epsilon(\bullet-z)\bigr)\right)_z
\cdot f(z)\,dz
=
T_{J\bar K}
\left(
\int_{\mathbb{C}^m}f(z)h_\epsilon(\bullet-z)\,dz
\right)
=
T_{J\bar K}(f*h_\epsilon).
\end{equation*}
この$T_\epsilon$ を$T$のsmoothingという. 
\end{defn}
\end{tcolorbox}


% [page 2]
%\noindent\underline{\textbf{Rem 13}}
\begin{rem}
\label{rem-H-5.13}
$H_{J\bar K,\epsilon} (z):=T_{J\bar K}\bigl(h_\epsilon(\bullet-z)\bigr)$とすると, これは$z$を変数とする$H_{J\bar K,\epsilon}\in C^\infty$であるので, 
$T_\epsilon$は$(p,q)$-formである. 
さらに$T_{J\bar K}*h_\epsilon=[H_{J\bar K,\epsilon}]$である. 
また $\sigma_{JK}\in\{\pm1\}$であって, 任意の$J\in\{m;p\}$, $K\in\{m;q\}$について
\begin{equation*}
\bigwedge_{\ell=1}^m
\left(dz^\ell\wedge d\bar z^\ell\right)
=
\sigma_{JK}\,
dz^J\wedge d\bar z^K\wedge
dz^{J^\circ}\wedge d\bar z^{K^\circ}
\end{equation*}
となるようにすると, 
任意の$\phi\in
\sum_{\substack{J\in\{m;m-p\}\\K\in\{m;m-q\}}}
\phi_{J\bar K}\,dz^J\wedge d\bar z^K
\in\mathcal{D}^{(m-p,m-q)}(U_\epsilon)$
について, 
\begin{align*}
T_\epsilon(\phi)
&=
\sum_{\substack{J\in\{m;p\}\\K\in\{m;q\}}}
\sigma_{JK}(T_{J\bar K}*h_\epsilon)
\left(\phi_{J^\circ\overline{K^\circ}}\right)
=
\sum_{\substack{J\in\{m;p\}\\K\in\{m;q\}}}
\sigma_{JK}
\int_{U_\epsilon}
H_{J\bar K}\cdot
\phi_{J^\circ\overline{K^\circ}}\,dz\\
&=
\int_{U_\epsilon}
\sum_{\substack{J\in\{m;p\}\\K\in\{m;q\}}}
\sigma_{JK}H_{J\bar K}\cdot
\phi_{J^\circ\overline{K^\circ}}\,dz
=
\int_{U_\epsilon}T_\epsilon\wedge\phi.
\end{align*}

特に$T_\epsilon$は
$
\left[
\sum_{\substack{J\in\{m;p\}\\K\in\{m;q\}}}
H_{J\bar K,\epsilon}\,dz^J\wedge d\bar z^K
\right].
$というsmooth $(p,q)$-formに一致する. 
\end{rem}



\begin{tcolorbox}[mybox]
\begin{lem}
\label{lem-H-5.14}
%\noindent\underline{\textbf{Lem 14}}
$U\subset\mathbb{C}^m$を開集合, 
$T\in\mathcal{D}'^{(p,p)}(U)$をcurrentとする. 
\begin{enumerate}
\item  $T$がpositiveならば, $\epsilon>0$で$T_\epsilon$もpositive.
\item  $\epsilon>0$を固定する. もし $dT_\epsilon=0$ がカレントとして成り立つならば, $dT_\epsilon=0$がformとして成り立つ. 
\end{enumerate}
\end{lem}
\end{tcolorbox}

\begin{proof}

(1). 任意の$f\in\mathcal{D}^{(m-p,m-p)}(U)$strongly positiveについて, 
$f \ast h_\epsilon$もstrongly positive. 
よって
$T_\epsilon(f)=T(f \ast h_\epsilon)\geq 0$であり, 
$T_\epsilon$はpositveである. 

(2). Lemma \ref{lem-H-5.12}から$dT_\epsilon=0$であり, 
\begin{equation*}
\int_{U_\epsilon}dT_\epsilon\wedge\phi=0
\qquad
(\forall\phi\in\mathcal{D}(U_\epsilon))
\end{equation*}
今
$
dT_\epsilon
=
\sum_{\substack{J,K\subset\{1,\ldots,m\}\\|J|+|K|=p+1}}
T_{\epsilon J\bar K}\,dz^J\wedge d\bar z^K
$なので, 
\begin{equation*}
\int_{U_\epsilon}T_{\epsilon J\bar K}\,\phi\,dz=0
\qquad
\text{for }\forall\phi\in\mathcal{D}(U_\epsilon).
\end{equation*}

Approximationとのたたみ込みを考えれば
$
T_{\epsilon J\bar K}=0
$が任意の$J,K$
で言えて, 
$dT_\epsilon=0$が$(p,p)$-formとして言える. 
\end{proof}


\begin{tcolorbox}[mybox]
\begin{lem}
\label{lem-H-5.15}
%\noindent\underline{\textbf{Lem 15}}
開集合$U\subset\mathbb{C}^m$, $T\in\mathcal{D}'^{(p,p)}(U)$をpositive currentとする. 
補足\ref{rem-H-5.8}によって, 
ある$\mu_{T_\epsilon\wedge\alpha^{m-p}}$というpositive real Radon measureがあって, 
\begin{equation*}
(T_\epsilon\wedge\alpha^{m-p})(f)
=
\int_U f\,d\mu_{T_\epsilon\wedge\alpha^{m-p}}
\qquad
\left(
\forall f\in\mathcal{K}(U_\epsilon)
\right).
\end{equation*}
となるものがある. この時任意のBorel可測集合$E\in\mathcal{B}(U_\epsilon)$について, 
\begin{equation*}
\mu_{T_\epsilon\wedge\alpha^{m-p}}(E)
=
\int_E T_\epsilon\wedge\alpha^{m-p}
\end{equation*}
\end{lem}
\end{tcolorbox}

\begin{proof}
$\alpha$の定義より, 
$
T_\epsilon\wedge\alpha^{m-p}
=
(T\wedge\alpha^{m-p})_\epsilon.
$
特に$T_\epsilon\wedge\alpha^{m-p}$
は$U_\epsilon$上のsmooth positive $(m,m)$-formである.
したがって, ある$C^\infty$関数
$g\colon U_\epsilon\longrightarrow\mathbb{R}$
かつ$
g(z)\geq0$で
\begin{equation*}
T_\epsilon\wedge\alpha^{m-p}
=
g
\left(
\bigwedge_{\ell=1}^m
\left(i\,dz^\ell\wedge d\bar z^\ell\right)
\right)
\end{equation*}
と書ける.
ここで
$
dV:=
\bigwedge_{\ell=1}^m
\left(i\,dz^\ell\wedge d\bar z^\ell\right)
$
とおく.
すると
$
T_\epsilon\wedge\alpha^{m-p}=g\,dV
$
である.

したがって, 任意の
$E\in\mathcal{B}(U_\epsilon)$に対して
\begin{equation*}
\int_E T_\epsilon\wedge\alpha^{m-p}
:=
\int_E g\,dV
\in[0,\infty]
\end{equation*}
と定義することができる.
ここで右辺はLebesgue積分である.

\noindent$\bullet$\;
$\nu\colon\mathcal{B}(U_\epsilon)\longrightarrow[0,\infty]$を
\begin{equation*}
\nu(E)
:=
\int_E T_\epsilon\wedge\alpha^{m-p}
=
\int_E g\,dV
\qquad
(E\in\mathcal{B}(U_\epsilon))
\end{equation*}
で定める.すると$\nu$はpositive real Radon measureになる(ここは議論が必要だが言える)
よって$\mu=\mu_{T_\epsilon\wedge\alpha^{m-p}}$の定義より,
任意の
$\phi\in \mathcal{K}^\infty(U_\epsilon)$に対して
\begin{align*}
\int_{U_\epsilon}\phi\,d\nu
=
\int_{U_\epsilon}
\phi\,T_\epsilon\wedge\alpha^{m-p}
=
(T_\epsilon\wedge\alpha^{m-p})(\phi)
=
\int_{U_\epsilon}\phi\,d\mu.
\end{align*}
よってRiesz表示定理の一意性より
$\mu=\nu$である. 
\end{proof}



\begin{tcolorbox}[mybox]
\begin{lem}
\label{lem-H-5.16}
% [page 1]
%\noindent\underline{\textbf{Lem 16}}
開集合$U\subset\mathbb{C}^m$, $T\in\mathcal{D}'^{(p,p)}(U)$ positive  currentとする. 
リースの表現定理より, ある
$\mu_{\operatorname{Tr}(T)}$となるpositive (real) Radon measureがあって, 
\begin{equation*}
(\operatorname{Tr}(T))(f)
=
\int_U f\,d\mu_{\operatorname{Tr}(T)}
\qquad
(\forall f\in\mathcal{K}(U)).
\end{equation*}
となる. 今, 固定された$z\in U$と$r>0$ で
$\overline{B(z;r)}\subset U$かつ$\mu_{\operatorname{Tr}(T)}(\Gamma(z;r))=0$となるならば, 
\begin{equation*}
\lim_{\epsilon\to0}
\int_{B(z;r)}T_\epsilon\wedge\alpha^{m-p}
=
\mu_{T\wedge\alpha^{m-p}}(B(z;r)).
\end{equation*}
%\textcolor{blue}{$T_\epsilon\wedge\alpha^{m-p}$は sm.\ form である}
\end{lem}
\end{tcolorbox}
\begin{proof}

[Step 1.] $\mu_{T\wedge\alpha^{m-p}}(\Gamma(z;r))=0$を示す. 
今
\begin{equation*}
T
=
i^{p^2}
\sum_{J,K\in\{m;p\}}
T_{J\bar K}\wedge dz^J\wedge d\bar z^K
\qquad
\text{here}\qquad
\operatorname{Tr}(T)
=
\sum_{J\in\{m;p\}}T_{J\bar J}.
\end{equation*}
とすると簡単な計算から
\begin{align*}
T\wedge\alpha^{m-p}
&=
T\wedge
\left(
\bigwedge^{m-p}
\frac{i}{2\pi}
\sum_{j=1}^m dz^j\wedge d\bar z^j
\right)\\
&=
\left(\frac{i}{2\pi}\right)^{m-p}
T\wedge
\left(
(m-p)!
\sum_{J\in\{m;m-p\}}
\underbrace{
dz^{j_1}\wedge d\bar z^{j_1}
\wedge\cdots\wedge
dz^{j_{m-p}}\wedge d\bar z^{j_{m-p}}
}_{=:(dz\wedge d\bar z)^J}
\right)\\
&=
(m-p)!
\left(\frac{i}{2\pi}\right)^{m-p}
i^{p^2}
\sum_{J\in\{m;m-p\}}
T_{J^\circ\overline{J^\circ}}\,
\underbrace{dz^{J^\circ}\wedge d\bar z^{J^\circ}}_{=
(-1)^{\frac12p(p-1)}(dz\wedge d\bar z)^{J^\circ}}
\wedge(dz\wedge d\bar z)^J.
\\
&=
(m-p)!
\frac{i^{m-p}}{(2\pi)^{m-p}}
i^{p^2}
\sum_{J\in\{m;m-p\}}
i^{p(p-1)}
T_{J^\circ\overline{J^\circ}}
(dz\wedge d\bar z)^{J^\circ}
\wedge(dz\wedge d\bar z)^J\\
&=
(m-p)!
\frac{i^m}{(2\pi)^{m-p}}
\sum_{J\in\{m;m-p\}}
T_{J^\circ\overline{J^\circ}}\,
dz^1\wedge d\bar z^1\wedge\cdots\wedge dz^m\wedge d\bar z^m\\
&=
\frac{(m-p)!}{(2\pi)^{m-p}}
\operatorname{Tr}(T)\wedge
\left(
\bigwedge_{\ell=1}^m
i\,dz^\ell\wedge d\bar z^\ell
\right).
\end{align*}
ここで$2p(p-1)\equiv0\pmod 4$を用いた. よって
\begin{equation*}
\frac{(m-p)!}{(2\pi)^{m-p}}
(\operatorname{Tr}(T))(f)
=
(T\wedge\alpha^{m-p})(f)
\qquad
(\forall f\in\mathcal{K}(U)).
\end{equation*}
リースの表現定理から, 
\begin{equation*}
\frac{(m-p)!}{(2\pi)^{m-p}}
\mu_{\operatorname{Tr}(T)}
=
\mu_{T\wedge\alpha^{m-p}}.
\end{equation*}
であるので, 
\begin{equation*}
\mu_{T\wedge\alpha^{m-p}}(\Gamma(z;r))
=
\frac{(m-p)!}{(2\pi)^{m-p}}
\mu_{\operatorname{Tr}(T)}(\Gamma(z;r))
=
0.
\end{equation*}

[Step 2] 任意の$\delta\in\mathbb{R}_{>0}$について, 
 Radon測度のregularityからある
$\sigma\in\mathbb{R}_{>0}$があって, 
\begin{equation}
\label{eq-support-delta}
0
\leq
\mu_{T\wedge\alpha^{m-p}}
\left(
B(z;r+3\sigma)\setminus\overline{B(z;r-3\sigma)}
\right)
<
\frac{1}{3}\delta.
\end{equation}
そこでbump関数($C^\infty$級関数)$f\colon\mathbb{R}\longrightarrow\mathbb{R}$で
\begin{equation*}
\begin{cases}
f(t)=1 & (\forall |t|\leq r-2\sigma),\\
0\leq f(t)\leq1 & (r-2\sigma\leq|t|\leq r-\sigma),\\
f(t)=0 & (|t|\geq r-\sigma).
\end{cases}
\end{equation*}
となるものを取り, $g(t):=1-f(t)$とする. 
$\bullet$を変数とみて関数を考えると, 
\begin{align*}
&
\left|
\int_{B(z;r)}T_\epsilon\wedge\alpha^{m-p}
-
\mu_{T\wedge\alpha^{m-p}}(B(z;r))
\right|\\
&=
\left|
\int_{U_\epsilon}
(f+g)(|\bullet-z|)\chi_{B(z;r)}(\bullet)
T_\epsilon\wedge\alpha^{m-p}
-
\int_U
(f+g)(|\bullet-z|)\chi_{B(z;r)}(\bullet)
\,d\mu_{T\wedge\alpha^{m-p}}
\right|\\
&\leq
\underbrace{
\left|
\int_{U_\epsilon}
f(|\bullet-z|)T_\epsilon\wedge\alpha^{m-p}
-
\int_U f(|\bullet-z|)
\,d\mu_{T\wedge\alpha^{m-p}}
\right|
}_{(3)}\\
&\quad+
\underbrace{\left|
\int_{B(z;r)}
g(|\bullet-z|)T_\epsilon\wedge\alpha^{m-p}
\right|}_{(2)}
+
\underbrace{\left|
\int_U
g(|\bullet-z|)\chi_{B(z;r)}(\bullet)
\,d\mu_{T\wedge\alpha^{m-p}}
\right|
}_{(1)}.
\end{align*}

(1)については
$\operatorname{Supp}\bigl(g(|\bullet-z|)\chi_{B(z;r)}(\bullet)\bigr)
\subset
B(z;r)\setminus\overline{B(z;r-2\sigma)}$
に注目して
\begin{align*}
0
&\leq
(1)
=
\int_U
g(|u-z|)\chi_{B(z;r)}(\bullet)
\,d\mu_{T\wedge\alpha^{m-p}}
\leq
\mu_{T\wedge\alpha^{m-p}}
\left(
B(z;r+3\sigma)\setminus\overline{B(z;r-3\sigma)}
\right)
\underalign{\eqref{eq-support-delta}}{<}
\frac{1}{3}\delta.
\tag{(*)}
\end{align*}

\begin{align*}
0
&\underalign{}{\leq}
(2) = \int_{B(z;r)}
g(|\bullet-z|)T_\epsilon\wedge\alpha^{m-p}
\underalign{}{=}
\int_U
g(|u-z|)\chi_{B(z;r)}(\bullet)
T_\epsilon\wedge\alpha^{m-p}
\\ 
&\underalign{\text{(Lem. \ref{lem-H-5.15})}}{=}
(T_\epsilon\wedge\alpha^{m-p})
\left(
g(|\bullet-z|)\chi_{B(z;r)}(\bullet)
\right)
\underalign{\text{(Def. \ref{defn-H-5.12})}}{=}
(T\wedge\alpha^{m-p})_\epsilon
\left(
g(|\bullet-z|)\chi_{B(z;r)}(\bullet)
\right)
\\
%&\underalign{\xb{Obs8}}{=}
%(T\wedge\alpha^{m-p})
%\left(
%g(|\bullet-z|)\chi_{B(z;r)}(\bullet)*h_\epsilon(\bullet)
%\right)\\
&\underalign{\text{(Rem. \ref{rem-H-5.8})}}{=}
\int_U
\underbrace{
g(|\bullet-z|)\chi_{B(z;r)}(\bullet)*h_\epsilon(\bullet)}_{\le \chi_{B(z;r+3\sigma)\setminus\overline{B(z;r-3\sigma)}} \text{ by Support \& $g, \chi \le 1$}}
\,d\mu_{T\wedge\alpha^{m-p}}
\underalign{}{\leq}
\int_U
\chi_{B(z;r+3\sigma)\setminus\overline{B(z;r-3\sigma)}}
\,d\mu_{T\wedge\alpha^{m-p}}
\\
&\underalign{}{=}
\mu_{T\wedge\alpha^{m-p}}
\left(
B(z;r+3\sigma)\setminus\overline{B(z;r-3\sigma)}
\right)
<
\frac{1}{3}\delta.
\tag{***}
\end{align*}

また$T_\epsilon\longrightarrow T \in \mathcal{D}'^{(p,p)}(U)$
が
$\epsilon\longrightarrow0$で言えるので, 
\cite[Theorem 6.32(b)]{Rud}から
ある$\epsilon_0\in(0,\epsilon)$であって, 
任意の$0<\epsilon'<\epsilon_0$について
\begin{align*}
&
|(3)|=
\left|
\int_{B(z;r)}
f(|\bullet-z|)T_{\epsilon'}\wedge\alpha^{m-p}
-
(T\wedge\alpha^{m-p})(f(|\bullet-z|))
\right|
=
\left|
(T\wedge\alpha^{m-p})
\left(
(f*h_{\epsilon'}-f)(|\bullet-z|)
\right)
\right|
<
\frac{\delta}{3}
\end{align*}
以上より
\begin{equation*}
|(3)|
+
|(2)|
+
|(1)|
<
\frac{\delta}{3}
+
\frac{\delta}{3}
+
\frac{\delta}{3}
=
\delta.
\end{equation*}
これは
\begin{equation*}
\lim_{\epsilon\to0}
\int_{B(z;r)}
T_\epsilon\wedge\alpha^{m-p}
=
\mu_{T\wedge\alpha^{m-p}}(B(z;r)).
\end{equation*}
を意味する. 
\end{proof}


\begin{tcolorbox}[mybox]
\begin{lem}
\label{lem-H-5.17}
%\noindent\underline{\textbf{Lem 17}}
$\beta=dd^c\log|z|$は $\mathbb{C}^m\setminus\{0\}$において$(1,1)$ positive currentである. 
\end{lem}
\end{tcolorbox}
\begin{proof}
\begin{equation*}
\beta
:=
dd^c\log|z|
=
\frac{i}{2\pi}
\sum_{j,k=1}^n
\left(
\frac{\delta_{jk}}{|z|^2}
-
\frac{z^k\bar z^j}{|z|^4}
\right)
dz^j\wedge d\bar z^k,
\end{equation*}
なので, 
\begin{equation*}
A
:=
\left(
|z|^2\delta_{jk}-z^k\bar z^j
\right)_{\substack{1\leq j\leq m\\1\leq k\leq m}}
=
|z|^2E_n
-
\begin{pmatrix}
\bar z_1\\
\vdots\\
\bar z_m
\end{pmatrix}
(z_1,\ldots,z_m)
\end{equation*}
がsemipositiveなHermitian 行列であること示せば良い. 
これは任意の$u:=(u_1,\ldots,u_m)\in\mathbb{C}^m$について, Cauchy--Schwarzの不等式を用いて
\begin{align*}
(u_1,\ldots,u_m)
A
\begin{pmatrix}
\bar u_1\\
\vdots\\
\bar u_m
\end{pmatrix}
&=
|u|^2|z|^2
-
(u_1,\ldots,u_m)
\begin{pmatrix}
\bar z_1\\
\vdots\\
\bar z_m
\end{pmatrix}
(z_1,\ldots,z_m)
\begin{pmatrix}
\bar u_1\\
\vdots\\
\bar u_m
\end{pmatrix}\\
&=
|u|^2|z|^2
-
|\langle u,z\rangle|^2
\geq0
\end{align*}
\end{proof}

\begin{tcolorbox}[mybox]
\begin{lem}
\label{lem-H-5.18}
%\noindent\underline{\textbf{Lem 18}}
開集合$(U\subset\mathbb{C}^m$, positive current $T\in\mathcal{D}'^{(p,p)}(U)$, 
$\operatorname{Tr}(T)$, $d\mu_{\operatorname{Tr}(T)}$を補題\ref{lem-H-5.16}の通りとする. 
$z\in U$, $r_2>r_1>0$で
$\overline{B(z;r_2)}\subset U$かつ
$\mu_{\operatorname{Tr}(T)}(\Gamma(z;r_1))
=
\mu_{\operatorname{Tr}(T)}(\Gamma(z;r_2))
=
0$であるとする. 
この時
\begin{equation*}
\lim_{\epsilon\to0}
\int_{B(z;r_2)\setminus\overline{B(z;r_1)}}
T_\epsilon\wedge\beta^{m-p}
=
\mu_{T\wedge\beta^{m-p}}
\left(
B(z;r_2)\setminus\overline{B(z;r_1)}
\right).
\end{equation*}
\end{lem}
\end{tcolorbox}
詳細は補題\ref{lem-H-5.16}とほぼ一緒なので省略.
%%%%%%%%%%%%%%%%%%
\begin{comment}


\xr{2026/08/17 ここからまとめる}
\begin{proof}
詳細は\xb{Lem 16}とほぼ一緒なのでアイデアのみのべる


まず，
\begin{equation*}
\mu_{T\wedge\beta^{m-p}}(\Gamma(z;r_1))
=
\mu_{T\wedge\beta^{m-p}}(\Gamma(z;r_2))
=
0
\end{equation*}
を示す．

（$T\wedge\beta^{m-p}$は$\operatorname{Tr}(T)$の定数倍で上から押さえられる．）

次に，$\forall\delta>0$に対し，
\begin{align*}
0
&\leq
\mu_{T\wedge\beta^{m-p}}
\Bigl(
\bigl(B(z;r_1+3\sigma)\setminus\overline{B(z;r_1-3\sigma)}\bigr)
\sqcup
\bigl(B(z;r_2+3\sigma)\setminus\overline{B(z;r_2-3\sigma)}\bigr)
\Bigr)\\
&<
\frac{1}{5}\delta
\qquad\cdots(*)
\end{align*}
となる$\sigma$をとり，$f\colon\mathbb{R}\longrightarrow\mathbb{R}$: $C^\infty$-func.\ s.t.
\begin{equation*}
\begin{cases}
f(t)=1
&
(r_1+2\sigma\leq|t|\leq r_2-2\sigma),\\
0\leq f(t)\leq1
&
(r_1+\sigma\leq|t|\leq r_1+2\sigma,\;
r_2-2\sigma\leq|t|\leq r_2-\sigma),\\
f(t)=0
&
(|t|\leq r_1+\sigma,\;
r_2+2\sigma\leq|t|)
\end{cases}
\end{equation*}
を選ぶ．

\textcolor{red}{[図：page01-fig1.png をここに挿入]}

% [page 2]
\begin{equation*}
g_1(t):=
\begin{cases}
1-f(t) & (|t|\leq r_1+2\sigma),\\
0 & (|t|\geq r_1+2\sigma),
\end{cases}
\qquad
g_2(t):=
\begin{cases}
0 & (|t|\leq r_2-2\sigma),\\
1-f(t) & (r_2-2\sigma\leq|t|).
\end{cases}
\end{equation*}

ととると，$f+g_1+g_2=1$．

あとは
\begin{align*}
&
\left|
\int_{B(z;r_2)\setminus\overline{B(z;r_1)}}
T_\epsilon\wedge\beta^{m-p}
-
\mu_{T\wedge\beta^{m-p}}
\left(
B(z;r_2)\setminus\overline{B(z;r_1)}
\right)
\right|\\
&=
\left|
\int_{U_\epsilon}
(f+g_1+g_2)(|u-z|)
\chi_{B(z;r_2)\setminus\overline{B(z;r_1)}}(u)
T_\epsilon\wedge\beta^{m-p}
\right.\\
&\hspace{45mm}\left.
-
\int_U
(f+g_1+g_2)(|u-z|)
\chi_{B(z;r_2)\setminus\overline{B(z;r_1)}}(u)
\,d\mu_{T\wedge\beta^{m-p}}
\right|\\
&\leq
\left|
\int_{U_\epsilon}
f(|u-z|)T_\epsilon\wedge\beta^{m-p}
-
\int_U
f(|u-z|)\,d\mu_{T\wedge\beta^{m-p}}
\right|
\quad\textcolor{blue}{\textcircled{5}}\\
&\quad+
\left|
\int_{B(z;r_2)\setminus\overline{B(z;r_1)}}
g_1(|u-z|)T_\epsilon\wedge\beta^{m-p}
\right|
\quad\textcolor{blue}{\textcircled{3}}\\
&\quad+
\left|
\int_{B(z;r_2)}
g_2(|u-z|)T_\epsilon\wedge\beta^{m-p}
\right|
\quad\textcolor{blue}{\textcircled{4}}\\
&\quad+
\left|
\int_U
g_1(|u-z|)
\chi_{B(z;r_2)\setminus\overline{B(z;r_1)}}(u)
\,d\mu_{T\wedge\beta^{m-p}}
\right|
\quad\textcolor{blue}{\textcircled{1}}\\
&\quad+
\left|
\int_U
g_2(|u-z|)
\chi_{B(z;r_2)}(u)
\,d\mu_{T\wedge\beta^{m-p}}
\right|
\quad\textcolor{blue}{\textcircled{2}}
\end{align*}
を全て$\dfrac{\delta}{5}$で上から押さえる．

\textcircled{1}と\textcircled{2}は$(*)$より簡単．\textcircled{5}はLem 16\textcircled{3}と同じ．\textcircled{3}と\textcircled{4}については
\begin{equation*}
C\alpha-\beta:\text{ pos.\ current on }
\overline{B(z;r_2+3\sigma)}\setminus B(z;r_1-3\sigma)
\end{equation*}
となる$C$が存在するので，$\beta$を$C\alpha$におきかえてLem 16の\textcircled{2}と同じ議論に帰着．\hfill //

\end{proof}
\xr{2026/08/17 ここからまとめる}
\end{comment}
%%%%%%%%%%%%%%%%%%%%%%
\begin{proof}[Proof of Theorem \ref{thm-H-5.10}]
$z=0$として良い. 
$dT=0$なので, カレントとして$dT_\epsilon=0$である. 
補題 \ref{lem-H-5.14}から$dT_\epsilon=0$は$C^\infty$級の$(2p+1)$-formである. 
補足 \ref{rem-H-5.4}の(Poincaré lemma)から
$C^\infty$級の$(2p-1)$-form$S_\epsilon$があって, $dS_\epsilon=T_\epsilon$である. 
また
\begin{equation*}
n(0;r,T_\epsilon)
\underalign{\text{(Def. \ref{defn-H-5.9})}}{=}
\frac{1}{r^{2(m-p)}}
\mu_{T_\epsilon\wedge\alpha^{m-1}}(B(0;r))
\underalign{\text{(Rem. \ref{rem-H-5.13})}}{=}
\frac{1}{r^{2(m-p)}}
\int_{B(0;r)}
T_\epsilon\wedge\alpha^{m-p}.
\end{equation*}

よって任意の$0< r_1\leq r_2$について, 
\begin{align*}
n(0;r_2,T_\epsilon)-n(0;r_1,T_\epsilon)
&\underalign{\text{}}{=}
\frac{1}{r_2^{2(m-p)}}
\int_{B(0;r_2)}
T_\epsilon\wedge\alpha^{m-p}
-
\frac{1}{r_1^{2(m-p)}}
\int_{B(0;r_1)}
T_\epsilon\wedge\alpha^{m-p}\\
&\underalign{\text{}}{=}
\frac{1}{r_2^{2(m-p)}}
\int_{B(0;r_2)}
dS_\epsilon\wedge\alpha^{m-p}
-
\frac{1}{r_1^{2(m-p)}}
\int_{B(0;r_1)}
dS_\epsilon\wedge\alpha^{m-p}\\
&\underalign{\text{(Stokes)}}{=}
\frac{1}{r_2^{2(m-p)}}
\int_{\Gamma(0;r_2)}
S_\epsilon\wedge\alpha^{m-p}
-
\frac{1}{r_1^{2(m-p)}}
\int_{\Gamma(0;r_1)}
S_\epsilon\wedge\alpha^{m-p}\\
&\underalign{\text{(Lem \ref{lem-H-5.3})}}{=}
\int_{\Gamma(0;r_2)}
S_\epsilon\wedge\beta^{m-p}
-
\int_{\Gamma(0;r_1)}
S_\epsilon\wedge\beta^{m-p}\\
&\underalign{\text{(Stokes)}}{=}
\int_{B(0;r_2)\setminus\overline{B(0;r_1)}}
T_\epsilon\wedge\beta^{m-p}.
\end{align*}
今$\widetilde r>r_2$で$\overline{B(0;\widetilde r)}\subset U$となるものをとる. 
任意の$N\in\mathbb{Z}_{>0}$について, 
\begin{equation*}
E_N:=
\left\{
r\in[0,\widetilde r)
\ \middle|\
\mu_{\operatorname{Tr}(T)}(\Gamma(r))>\frac{1}{N}
\right\}.
\end{equation*}
とすると以下の評価が成り立つ: 
\begin{equation*}
\infty
>
\mu_{\operatorname{Tr}(T)}
\left(
\overline{B(0;\widetilde r)}
\right)
\geq
\sum_{r'\in E_N}
\mu_{\operatorname{Tr}(T)}(\Gamma(0;r'))
\geq
\#E_N\cdot\frac{1}{N}.
\end{equation*}
特に$E_N$は有限集合である.
よって$E:=\{r\in[0,\widetilde r) \mid \mu_{\operatorname{Tr}(T)}(\Gamma(0;r))>0\}$は可算集合である. 
今$r_1,r_2\notin E$ならば
\begin{equation*}
\mu_{\operatorname{Tr}(T)}(\Gamma(0;r_1))
=
\mu_{\operatorname{Tr}(T)}(\Gamma(0;r_2))
=
0,
\end{equation*}
なので
\begin{align*}
n(0;r_2,T)-n(0;r_1,T)
&\underalign{\text{(Lem. \ref{lem-H-5.16})}}{=}
\lim_{\epsilon\to0}
\left(
n(0;r_2,T_\epsilon)-n(0;r_1,T_\epsilon)
\right)\\
&\underalign{\text{}}{=}
\lim_{\epsilon\to0}
\int_{B(0;r_2)\setminus\overline{B(0;r_1)}}
T_\epsilon\wedge\beta^{m-p}
\underalign{\text{(Lem. \ref{lem-H-5.18})}}{=}
\mu_{T\wedge\beta^{m-p}}
\left(
B(z;r_2)\setminus\overline{B(z;r_1)}
\right)
\geq0.
\end{align*}
また正則性から$n(0;r,T)$は$r$に関して右連続である. よって$r$に関して非減少関数である. 
\end{proof}
実はLemma \ref{lem-H-5.18}はなくても証明はできている(0以上はわかる.)

\begin{tcolorbox}[mybox]
\begin{defn}[{Lelong number}]
\label{defn-H-5.19}
%\noindent\underline{\textbf{Def 1.19 (Lelong number)}}
$T\in\mathcal{D}'^{(p,p)}(U)$をpositive current で$dT=0$となるものとする. 
$z\in U$について, Lelong number $\mathcal{L}(z;T)$を以下のように定義する. 
\begin{equation*}
\mathcal{L}(z;T)
:=
\lim_{r\to0}
\frac{1}{r^{2(m-p)}}
\mu_{T\wedge\alpha^{m-p}}(B(z;r))
\in\mathbb{R}_{\geq0}.
\end{equation*}
\end{defn}
\end{tcolorbox}



% [PDF page 2]
\section{Lelong number for complex submanifolds}
\label{sec-H-0925}

\underline{Notation（前回より）}
\begin{itemize}
\item $(z^1,\ldots,z^m)$: coordinate of $\C^m$.
\item $d^c:=\dfrac{1}{2\pi}\dfrac{i}{2}(\bar\partial-\partial)$
\quad (thus $dd^c=\dfrac{i}{2\pi}\partial\bar\partial$).
\item $\displaystyle
\alpha:=dd^c|z|^2
=\frac{i}{2\pi}\sum_{j=1}^m dz^j\wedge d\bar z^j$.
\item $\displaystyle
\beta:=dd^c\log|z|^2
=\frac{i}{2\pi}\sum_{j,\ell=1}^m
\left(\frac{\delta_{j\ell}}{|z|^2}
-\frac{z^\ell\bar z^j}{|z|^4}\right)
dz^j\wedge d\bar z^\ell$
\quad on $\C^m\setminus\{0\}$.
\item $\alpha^q$, $\beta^q$: $q$回のウェッジ積.
\quad $\alpha^0:=1$, $\beta^0:=1$.
\item $\eta:=d^c\log|z|^2\wedge\beta^{m-1}$.
\item $B(z;r):=\{w\in\C^m\mid |w-z|<r\}$.
\item $\Gamma(z;r):=\{w\in\C^m\mid |w-z|=r\}$.
\item $B(r):=B(0;r)$, $\Gamma(r):=\Gamma(0;r)$.
\end{itemize}

$U\subset\C^m$を開集合,
$S\subset U$を$k$次元のclosed complex submanifoldとする.
ここでclosedとは$U$の相対位相に関して閉じているという意味である.
$S$は積分カレント
$I_S\in\mathcal{D}'^{(m-k,m-k)}(U)$を
\begin{equation*}
I_S(\phi):=\int_S\phi
\qquad
\left(\phi\in\mathcal{D}^{(k,k)}(U)\right)
\end{equation*}
によって定める. $I_S$はpositive currentである.
また$S$は境界を持たないので, Stokesの定理から$dI_S=0$である.

$K\subset U$をコンパクト集合とする.
$\Vol(K\cap S)$は標準Euclid計量から誘導される体積である.
具体的には, $z^j=x^j+i\,y^j$として
\begin{equation*}
g:=\sum_{j=1}^m
\left(dx^j\otimes dx^j+dy^j\otimes dy^j\right)
\end{equation*}
とおく. $S$上の実局所座標$(u^1,\ldots,u^{2k})$を用いて
\begin{equation*}
\left.g\right|_S
=\sum_{i,j=1}^{2k}g_{ij}\,du^i\otimes du^j
\end{equation*}
と書くと, 複素構造が定める向きと整合的な座標について
\begin{equation*}
\mathrm{vol}_S
:=\sqrt{\det(g_{ij})_{1\leq i,j\leq 2k}}\,
 du^1\wedge\cdots\wedge du^{2k},
\qquad
\Vol(K\cap S):=\int_{K\cap S}\mathrm{vol}_S.
\end{equation*}

% [PDF page 3; Rem 1]
\begin{rem}
\label{rem-H-0925-volume}
上の記号のもとで,
\begin{equation*}
\Vol(K\cap S)=\frac{\pi^k}{k!}\int_{K\cap S}\alpha^k.
\end{equation*}
実際, 標準複素構造を$J$とすると,
$g$に対応する$2$-formは
\begin{equation*}
\omega:=g(J\,\cdot,\cdot)
=\sum_{j=1}^m dx^j\wedge dy^j
=\frac{i}{2}\sum_{j=1}^m dz^j\wedge d\bar z^j
=\pi\alpha.
\end{equation*}
適当な局所座標で計算すれば
\begin{equation*}
\mathrm{vol}_S
=\frac{1}{k!}\left(\left.\omega\right|_S\right)^k
=\frac{\pi^k}{k!}\left.\alpha^k\right|_S
\end{equation*}
となるので, 主張が従う.

特に$U=\C^m$, $S=V$を$k$次元複素線形部分空間,
$K=\overline{B(r)}$とすると,
\begin{align*}
\int_{B(r)\cap V}\alpha^k
&=\int_{\overline{B(r)}\cap V}\alpha^k\\
&=\frac{k!}{\pi^k}\Vol\left(\overline{B(r)}\cap V\right)
=r^{2k}.
\end{align*}
最後の等号では, 前回のFact 2にある$2k$次元Euclid球の体積
$\dfrac{\pi^k}{k!}r^{2k}$を用いた.
\end{rem}

% 原稿では Lem 3．記載順（Rem 1, Lem 3, Prop 2, Cor 3）はそのまま．
\begin{tcolorbox}[mybox]
\begin{lem}
\label{lem-H-0925-measure}
$\mu_{I_S\wedge\alpha^k}$をpositive current
$I_S\wedge\alpha^k$に対応するRadon測度とする（前回のObs 8）.
$z\in U$, $r>0$, $\overline{B(z;r)}\subset U$に対して,
前回のDef 9と同様に
\begin{equation*}
n(z;r,I_S):=\frac{1}{r^{2k}}
\mu_{I_S\wedge\alpha^k}(B(z;r))
\end{equation*}
と定める. このとき
\begin{equation*}
n(z;r,I_S)=\frac{1}{r^{2k}}\int_{B(z;r)\cap S}\alpha^k.
\end{equation*}
特に$n(r;I_S):=n(0;r,I_S)$と書く.
\end{lem}
\end{tcolorbox}

\begin{proof}
$S$上のBorel測度$\nu$を
\begin{equation*}
\nu(E):=\int_E\left(\left.\alpha\right|_S\right)^k
\qquad (E\in\mathcal{B}(S))
\end{equation*}
によって定める.
補足\ref{rem-H-0925-volume}より
$\left(\left.\alpha\right|_S\right)^k
=\dfrac{k!}{\pi^k}\mathrm{vol}_S$なので, $\nu$はRadon測度である.
包含写像$\iota\colon S\hookrightarrow U$による押し出しを
$\mu:=\iota_*\nu$とおく. すなわち
\begin{equation*}
\mu(E):=\nu(E\cap S)=\int_{E\cap S}\alpha^k
\qquad (E\in\mathcal{B}(U)).
\end{equation*}
$S$は$U$で閉じているので, $\mu$も$U$上のRadon測度である.
任意の$\phi\in\mathcal{D}(U)$について
\begin{align*}
(I_S\wedge\alpha^k)(\phi)
&=I_S(\alpha^k\wedge\phi)
=\int_S\phi\,\alpha^k
=\int_U\phi\,d\mu.
\end{align*}
よってRiesz表示定理の一意性から
$\mu=\mu_{I_S\wedge\alpha^k}$となる.
特に
\begin{equation*}
\mu_{I_S\wedge\alpha^k}(B(z;r))
=\int_{B(z;r)\cap S}\alpha^k
\end{equation*}
であり, 主張が従う.
\end{proof}
% 補足: 原稿ではS上の測度とU上の測度が同じ記号で書かれているため，
% 包含写像による押し出しを明記した．

% [PDF page 4; Prop 2]
\begin{tcolorbox}[mybox]
\begin{prop}[{野口--落合, (3.2.39)}]
\label{prop-H-0925-lelong}
$U\subset\C^m$を開集合, $S\subset U$を$k$次元の
closed complex submanifoldとする.
$z\in S$, $r>0$, $\overline{B(z;r)}\subset U$とする.
$I_S$の$z$におけるLelong numberを
\begin{equation*}
\mathcal{L}(z;I_S):=\lim_{t\to0+}n(z;t,I_S)
\end{equation*}
とおく（前回のDef 19）. このとき
\begin{equation*}
n(z;r,I_S)
=\frac{1}{r^{2k}}\int_{B(z;r)\cap S}\alpha^k
\geq\mathcal{L}(z;I_S)=1.
\end{equation*}
\end{prop}
\end{tcolorbox}
% 原稿で省略されている z\in S と球の条件を明記した．

\begin{proof}
平行移動によって$z=0$としてよい.
補題\ref{lem-H-0925-measure}と前回のThm 10
（closed positive currentの$n(z;r,T)$は$r$について非減少）から,
$\mathcal{L}(0;I_S)=1$を示せば十分である.
以下$1\leq k<m$の場合を示す.
$k=m$の場合は$S$が$0$の近傍で開集合となることから,
$k=0$の場合は十分小さい球内の$S$が$\{0\}$となることから従う.

[Step 1.] $S$を接空間上の正則グラフとして表示する.

$S$は$0$の近傍で$k$次元複素部分多様体なので,
$0$のある近傍$W\subset U$上の正則関数
$f_1,\ldots,f_{m-k}$があって,
\begin{equation*}
S\cap W=\{z\in W\mid f_1(z)=\cdots=f_{m-k}(z)=0\},
\qquad
\rank\left(\frac{\partial f_i}{\partial z^j}(0)\right)
_{\substack{1\leq i\leq m-k\\1\leq j\leq m}}=m-k.
\end{equation*}
ユニタリ座標変換を施すことにより,
標準Euclid計量と$\alpha$を保ったまま
\begin{equation*}
T_0S
=\left\langle\frac{\partial}{\partial z^1},\ldots,
\frac{\partial}{\partial z^k}\right\rangle
\subset T_0\C^m
\end{equation*}
としてよい. このとき
\begin{equation*}
\left(\frac{\partial f_i}{\partial z^j}(0)\right)
_{\substack{1\leq i\leq m-k\\1\leq j\leq k}}=0.
\end{equation*}

% [PDF page 5]
したがって
\begin{equation*}
\det\left(\frac{\partial f_i}{\partial z^{k+j}}(0)\right)
_{1\leq i,j\leq m-k}\neq0.
\end{equation*}
$f_1,\ldots,f_{m-k}$を可逆な定数行列による線形結合で取り替えれば,
\begin{equation*}
\left(\frac{\partial f_i}{\partial z^{k+j}}(0)\right)
_{1\leq i,j\leq m-k}=E_{m-k}
\end{equation*}
としてよい. ここで$E_{m-k}$は単位行列である.
% 原稿の「det(行列)=E_{m-k}」は「行列=E_{m-k}」に修正．

正則陰関数定理から, $0\in\C^k$の近傍で定義された正則関数
$F^{k+1},\ldots,F^m$が存在し, $W$を縮めれば
\begin{gather*}
F^{k+j}(0)=0\qquad(1\leq j\leq m-k),\\
S\cap W
=\left\{(z^1,\ldots,z^m)\in W\ \middle|\
z^{k+j}=F^{k+j}(z^1,\ldots,z^k),\ 1\leq j\leq m-k\right\}.
\end{gather*}
$w=(z^1,\ldots,z^k)$, $F=(F^{k+1},\ldots,F^m)$とおく.
$f_i(w,F(w))=0$を$w^j$で微分し, $w=0$で評価すると
\begin{align*}
0
&=\frac{\partial f_i}{\partial z^j}(0)
+\sum_{\ell=1}^{m-k}
\frac{\partial f_i}{\partial z^{k+\ell}}(0)
\frac{\partial F^{k+\ell}}{\partial w^j}(0)\\
&=0+\frac{\partial F^{k+i}}{\partial w^j}(0)
\qquad(1\leq i\leq m-k,\ 1\leq j\leq k).
\end{align*}
つまり$F(0)=0$, $dF(0)=0$である.
したがって, 各$F^{k+i}$の$0$の近傍でのべき級数展開には
$2$次以上の項しか現れない. 特に
\begin{equation*}
F(w)=O(|w|^2),\qquad dF(w)=O(|w|).
\end{equation*}

[Step 2.] 体積の主要項を計算する.

\xb{以下では, 原稿の積分領域の同一視を補足して記述する.}
\begin{equation*}
V:=\{(z^1,\ldots,z^k,0,\ldots,0)\in\C^m\}\simeq\C^k,
\qquad
\Phi(w):=(w,F(w))
\end{equation*}
とおき, $S$の$0$の近傍を$V$上のグラフで表示する.
$B(r)\cap S$に対応する$V$上の領域は, 厳密には$B(r)\cap V$ではなく
\begin{equation*}
D_r:=\{w\in\C^k\mid |w|^2+|F(w)|^2<r^2\}
\end{equation*}
である（$r>0$は十分小さくとる）.
$V$上の標準形式を
\begin{equation*}
\alpha_V:=\frac{i}{2\pi}\sum_{j=1}^k dw^j\wedge d\bar w^j
\end{equation*}
と書けば,
\begin{align*}
\Phi^*\alpha
&=\frac{i}{2\pi}\left(
\sum_{j=1}^k dw^j\wedge d\bar w^j
+\sum_{\ell=1}^{m-k}dF^{k+\ell}\wedge d\overline{F^{k+\ell}}
\right),\\
n(r;I_S)
&=\frac{1}{r^{2k}}\int_{D_r}(\Phi^*\alpha)^k.
\end{align*}
$dF(w)=O(|w|)$より, 体積形式の比として
\begin{equation*}
(\Phi^*\alpha)^k
=\left(1+O(|w|^2)\right)\alpha_V^k.
\end{equation*}
また$|F(w)|\leq C|w|^2$となる定数$C>0$をとれば,
十分小さい$r$について
\begin{equation*}
B_V\left(\frac{r}{\sqrt{1+C^2r^2}}\right)
\subset D_r\subset B_V(r),
\qquad B_V(t):=B(t)\cap V.
\end{equation*}
補足\ref{rem-H-0925-volume}より
$\displaystyle\int_{B_V(t)}\alpha_V^k=t^{2k}$なので,
\begin{equation*}
(1+C^2r^2)^{-k}
\leq\frac{1}{r^{2k}}\int_{D_r}\alpha_V^k\leq1.
\end{equation*}
以上から
\begin{equation*}
n(r;I_S)
=\frac{1}{r^{2k}}\int_{D_r}
\left(1+O(|w|^2)\right)\alpha_V^k
=1+O(r^2).
\end{equation*}

% [PDF page 6]
特に原稿にある評価$n(r;I_S)=1+O(r)$も従い,
\begin{equation*}
\mathcal{L}(0;I_S)=\lim_{r\to0+}n(r;I_S)=1.
\end{equation*}
これで命題が示された.
\end{proof}

% [PDF page 6; Cor 3]
\begin{tcolorbox}[mybox,breakable=false]
\begin{cor}[{野口--落合, (3.2.40)}]
\label{cor-H-0925-volume-lower}
命題\ref{prop-H-0925-lelong}と同じ仮定のもとで,
\begin{equation*}
\Vol(B(z;r)\cap S)\geq\frac{\pi^k}{k!}r^{2k}.
\end{equation*}
特に$0\in S$の場合,
$\Vol(B(r)\cap S)\geq\dfrac{\pi^k}{k!}r^{2k}$である.
\end{cor}
\end{tcolorbox}

\begin{proof}
命題\ref{prop-H-0925-lelong}から
\begin{equation*}
\frac{1}{r^{2k}}\int_{B(z;r)\cap S}\alpha^k\geq1,
\qquad\text{すなわち}\qquad
\int_{B(z;r)\cap S}\alpha^k\geq r^{2k}.
\end{equation*}
補足\ref{rem-H-0925-volume}の体積形式の等式を用いて,
\begin{equation*}
\Vol(B(z;r)\cap S)
=\frac{\pi^k}{k!}\int_{B(z;r)\cap S}\alpha^k
\geq\frac{\pi^k}{k!}r^{2k}.
\end{equation*}
\end{proof}

\begin{rem}[原稿の表記と補足について]
\label{rem-H-0925-editorial}
\leavevmode\par
\begin{itemize}
\item $\beta=dd^c\log|z|^2$は今回のPDFの表記に合わせた.
\item 命題・系には$z\in S$と$\overline{B(z;r)}\subset U$を明記した.
\item 原稿の$S$上の測度は, 包含写像によって$U$上に押し出して扱った.
\item 座標変換はEuclid計量を保つユニタリ変換とし,
原稿の「$\det(\text{行列})=E_{m-k}$」は「行列$=E_{m-k}$」に整えた.
\item グラフ表示による積分領域を$D_r$として明記し,
原稿の$1+O(r)$に至る議論を補った.
\end{itemize}
\end{rem}




